% ch8_b (书页 336-348 / p351-p363)
%--A-TAIL--
%JOIN%
% ==== tail：8.3 节（The Friedmann Equation）收尾，承接书页 335 末的 (8.64)，
% ==== 至书页 337 末止；置于本文件第一个 \section 之前 ====
% ---- 书页 336 / p351 ----
$\mu\nu=00$ 分量方程为
\begin{equation*}
-3\frac{\ddot{a}}{a} = 4\pi G(\rho + 3p),
\tag{8.65}
\end{equation*}
而 $\mu\nu=ij$ 各分量方程给出
\begin{equation*}
\frac{\ddot{a}}{a} + 2\left(\frac{\dot{a}}{a}\right)^{2} + 2\frac{\kappa}{a^{2}} = 4\pi G(\rho - p).
\tag{8.66}
\end{equation*}
由于各向同性，由 $\mu\nu=ij$ 只能得到一个独立的方程。我们可以用式 (8.65) 消去式 (8.66) 中的二阶导数，再稍加整理，得到
\begin{equation*}
\boxed{\,\left(\frac{\dot{a}}{a}\right)^{2} = \frac{8\pi G}{3}\rho - \frac{\kappa}{a^{2}},\,}
\tag{8.67}
\end{equation*}
以及
\begin{equation*}
\boxed{\,\frac{\ddot{a}}{a} = -\frac{4\pi G}{3}(\rho + 3p).\,}
\tag{8.68}
\end{equation*}
这两条方程合称\textbf{弗里德曼方程}（Friedmann equations），而具有式 (8.43) 的形式且满足这些方程的度规，定义了 Friedmann--Robertson--Walker（FRW）宇宙。事实上，如果我们知道 $\rho$ 对 $a$ 的依赖关系，其中第一条式 (8.67) 就足以解出 $a(t)$；当你听到人们提起\emph{那条}弗里德曼方程时，指的就是它，而式 (8.68) 有时被称为\emph{第二}弗里德曼方程。

有一大堆与宇宙学参数相联系的术语，我们在这里只介绍最基本的部分。膨胀的速率由\textbf{Hubble 参量}（Hubble parameter）表征，
\begin{equation*}
\boxed{\,H = \frac{\dot{a}}{a}.\,}
\tag{8.69}
\end{equation*}

% ---- 书页 337 / p352 ----
Hubble 参量在当前时期的值就是 Hubble 常数，$H_0$。目前的测量使我们相信，Hubble 常数为 $70\pm10$ km/sec/Mpc。（Mpc 代表兆秒差距（megaparsec），即 $3.09\times10^{24}$ cm。）由于这个数值仍存在一些不确定性，我们常把 Hubble 常数参数化为
\begin{equation*}
H_0 = 100\,h\ \mathrm{km/sec/Mpc},
\tag{8.70}
\end{equation*}
于是 $h\approx0.7$。典型的宇宙学尺度由\textbf{Hubble 长度}（Hubble length）设定
\begin{align*}
d_H &= H_0^{-1}c \\
&= 9.25\times10^{27}\,h^{-1}\ \mathrm{cm} \\
&= 3.00\times10^{3}\,h^{-1}\ \mathrm{Mpc},
\tag{8.71}
\end{align*}
以及\textbf{Hubble 时间}（Hubble time）
\begin{align*}
t_H &= H_0^{-1} \\
&= 3.09\times10^{17}\,h^{-1}\ \mathrm{sec} \\
&= 9.78\times10^{9}\,h^{-1}\ \mathrm{yr}.
\tag{8.72}
\end{align*}
当然，由于我们通常取 $c=1$，你会看到 $H_0^{-1}$ 既被称作 Hubble 长度，又被称作 Hubble 时间。此外还有\textbf{减速参数}（deceleration parameter），
\begin{equation*}
q = -\frac{a\ddot{a}}{\dot{a}^{2}},
\tag{8.73}
\end{equation*}
它量度膨胀速率的变化速率。

另一个有用的量是\textbf{密度参数}（density parameter），
\begin{equation*}
\boxed{\,\Omega = \frac{8\pi G}{3H^{2}}\rho = \frac{\rho}{\rho_{\mathrm{crit}}},\,}
\tag{8.74}
\end{equation*}
其中\textbf{临界密度}（critical density）定义为
\begin{equation*}
\rho_{\mathrm{crit}} = \frac{3H^{2}}{8\pi G}.
\tag{8.75}
\end{equation*}
这个量一般会随时间变化，被称为\emph{临界}密度，因为弗里德曼方程 (8.67) 可以写成
\begin{equation*}
\Omega - 1 = \frac{\kappa}{H^{2}a^{2}}.
\tag{8.76}
\end{equation*}
因此，$\kappa$ 的符号由 $\Omega$ 大于、等于还是小于 1 决定。我们有
\begin{align*}
\rho < \rho_{\mathrm{crit}} \quad&\leftrightarrow\quad \Omega < 1 \quad\leftrightarrow\quad \kappa < 0 \quad\leftrightarrow\quad \text{开放} \\
\rho = \rho_{\mathrm{crit}} \quad&\leftrightarrow\quad \Omega = 1 \quad\leftrightarrow\quad \kappa = 0 \quad\leftrightarrow\quad \text{平直} \\
\rho > \rho_{\mathrm{crit}} \quad&\leftrightarrow\quad \Omega > 1 \quad\leftrightarrow\quad \kappa > 0 \quad\leftrightarrow\quad \text{闭合}.
\end{align*}
于是，密度参数告诉我们，三种 Robertson--Walker 几何中哪一种描述我们的宇宙。从观测上确定它是至关重要的；近来对宇宙微波背景各向异性的测量使我们相信，$\Omega$ 非常接近于 1。

% ---- 书页 338 / p353 ----

\section{尺度因子的演化}

给定不同组分 $i$ 的能量密度 $\rho_i$ 的数值，连同它们的物态方程 $p_i=p_i(\rho_i)$，以及空间曲率 $\kappa$ 的大小，就可以求解弗里德曼方程 (8.67)，得到尺度因子 $a(t)$ 演化的完整历史。一般来说，我们只需对弗里德曼方程（它只是一阶微分方程）作数值积分，不过，对不同宇宙学参数所对应的解的类型获得一些直观感受是有用的。

为了简化任务，让我们设想能量密度的各个不同组分都按幂律演化，
\begin{equation*}
\rho_i = \rho_{i0}a^{-n_i}.
\tag{8.77}
\end{equation*}
与式 (8.55) 相比较，这等价于假定每个物态方程参数 $w_i=p_i/\rho_i$ 都是一个常数，其值为
\begin{equation*}
w_i = \frac{1}{3}n_i - 1.
\tag{8.78}
\end{equation*}
把空间曲率的贡献当作一种虚构的能量密度（fictitious energy density）来处理，我们还可以进一步简化表达式：
\begin{equation*}
\rho_{\mathrm{c}} \equiv -\frac{3\kappa}{8\pi G a^{2}},
\tag{8.79}
\end{equation*}
以及相应的密度参数
\begin{equation*}
\Omega_{\mathrm{c}} = -\frac{\kappa}{H^{2}a^{2}}.
\tag{8.80}
\end{equation*}
当然，它\emph{不是}能量密度，所以别忘了这只是一种记号上的障眼法。我们最喜欢的几种源的行为总结在下表中。
\begin{equation*}
\begin{array}{l|c|c}
 & w_i & n_i \\ \cline{2-3}
\text{物质} & 0 & 3 \\
\text{辐射} & \frac{1}{3} & 4 \\
\text{曲率} & -\frac{1}{3} & 2 \\
\text{真空} & -1 & 0
\end{array}
\tag{8.81}
\end{equation*}
在这些变量下，弗里德曼方程 (8.67) 可以写成
\begin{equation*}
H^{2} = \frac{8\pi G}{3}\sum_{i(\mathrm{c})}\rho_i,
\tag{8.82}
\end{equation*}
其中记号 $\sum_{i(\mathrm{c})}$ 表示我们不仅对所有真实的能量密度组分 $\rho_i$ 求和，也对空间曲率的贡献

% ---- 书页 339 / p354 ----
$\rho_{\mathrm{c}}$ 求和。注意，如果把两边除以 $H^{2}$，就得到
\begin{equation*}
1 = \sum_{i(\mathrm{c})}\Omega_i.
\tag{8.83}
\end{equation*}
右边的和\emph{不是}总密度参数 $\Omega$，后者只包含来自真实能量密度（不含曲率）的贡献；因此我们有
\begin{equation*}
\Omega_{\mathrm{c}} = 1 - \Omega.
\tag{8.84}
\end{equation*}

让我们先问一问：如果所有的 $\rho_i$（包括 $\rho_{\mathrm{c}}$）都非负，会发生什么。由于 $H^{2}$ 正比于 $\sum_{i(\mathrm{c})}\rho_i$，只要 $\sum_{i(\mathrm{c})}\rho_i\neq0$，宇宙就永远不会经历从膨胀到收缩的转变。我们还可以对 Hubble 参量求时间导数，
\begin{equation*}
\dot{H} = \frac{\ddot{a}}{a} - \left(\frac{\dot{a}}{a}\right)^{2},
\tag{8.85}
\end{equation*}
并代入两条弗里德曼方程 (8.67) 和 (8.68)，得到
\begin{equation*}
\dot{H} = -4\pi G\sum_{i(\mathrm{c})}(1+w_i)\rho_i.
\tag{8.86}
\end{equation*}
由于我们设想 $|w_i|\le1$，当所有 $\rho_i$ 都非负时，我们将总有 $\dot{H}\le0$。换言之，宇宙持续膨胀，但膨胀速率不断减小（这就引出一个绝佳的问题：是什么让宇宙一开始就如此之大？）。

由式 (8.85) 我们看到，$\ddot{a}$ 可以在 $\dot{H}$ 为负的同时为正——即使以 Hubble 参量量度的膨胀速率在减小，尺度因子也可以在“加速”（例如当 $a\propto t^{2}$ 时）。这是非欧几何不可避免的一个微妙之处。Hubble 参量与尺度因子的导数回答的是两个不同的问题。如果把两个检验粒子放在某个固定的初始距离上，问此后很短的时间内它们分开了多少，答案由 Hubble 参量给出；反之，如果选定某个固定的源，问它随时间看上去如何远离我们，答案则由尺度因子的变化给出。因此，“加速”（或“减速”）有两种截然不同但同样正当的含义。实际上，“加速”通常指 $\ddot{a}>0$ 的情形，即使 $\dot{H}<0$。这番讨论并非纯粹学究式的；正如我们在下文将要看到的，我们当前真实的宇宙似乎正是这种类型。

每个 $\rho_i$ 都非负绝非必然。物质和辐射来自动力学的粒子与场，因此我们预期它们的能量密度永远不会为负；如果它们可以为负，空的空间就可能衰变成一团正、负能量的场。但真空和曲率就是另一回事了。真空能是非动力学的，所以负值不会诱发任何不稳定性；而曲率只是空间几何的一种属性，正负皆可。于是，如果我们有
% ---- 书页 340 / p355 ----
负的真空能或正的空间曲率（记住 $\rho_{\mathrm{c}}\propto-\kappa$），Hubble 参量就可以变为零甚至变号。de Sitter 度规 (8.6) 就提供了一个例子：它既有正的真空能，又有正的空间曲率；它描述了一个先坍缩、到达转折点、然后开始膨胀的宇宙。

真实世界是个不整洁的地方，由许许多多不同种类的能量密度组成。不过，由于不同的源以不同的速率演化，在很长一段时间内能量密度会明显地由某一种源主导。因此，考察弗里德曼方程在只有一种能量密度 $\rho\propto a^{-n}$ 时的解是非常有用的。因为我们把空间曲率也当作一种有效的能量源包括进来，这意味着我们要么考虑由单一源主导的平直宇宙，要么考虑具有空间曲率的完全空无一物的宇宙。弗里德曼方程于是意味着
\begin{equation*}
\dot{a}\propto a^{1-n/2}.
\tag{8.87}
\end{equation*}
这可以立即积分得到
\begin{equation*}
\boxed{\,a\propto t^{2/n}\qquad(\text{当 }\rho\propto a^{-n}\text{ 时}).\,}
\tag{8.88}
\end{equation*}
例如考虑一个由物质主导的平直宇宙，$\Omega=\Omega_{\mathrm{M}}=1$；它被称为 Einstein--de Sitter 模型，并且在很长一段时间内曾是描述真实世界的最受青睐的模型（至少在理论家中间是如此）。在 Einstein--de Sitter 宇宙中，尺度因子按 $a\propto t^{2/3}$ 演化；而平直的辐射主导宇宙则按 $a\propto t^{1/2}$ 演化。任何 $n>2$ 的这类宇宙的共形图在附录 H 中导出。尽管我们相信真实宇宙中物质、辐射和真空能的数量都不为零，这些解仍然非常有用；正如我们后面将要讨论的，宇宙在早期是辐射主导的，而在宇宙从 $a\sim1/3000$ 膨胀到 $a\sim1/2$ 的过程中则是物质主导的。

这些解都有一个 $a=0$ 处的奇点，称为\textbf{大爆炸}（Big Bang）。它代表宇宙从奇性状态的创生，而不是物质爆炸进入一个先前存在的时空。人们或许会希望，我们的 FRW 宇宙的完美对称性正是这个奇点的起因，但事实并非如此；宇宙学奇点定理（cosmological singularity theorems）表明，任何满足 $\rho>0$ 且 $p\ge0$ 的宇宙都必定始于一个奇点。当然，当 $a\to0$ 时能量密度变得任意高，我们不指望经典广义相对论在这个区域还能准确地描述自然；想必量子引力将变得重要，尽管目前尚不清楚它如何起作用。

看式 (8.88) 可见，真空能主导（$n=0$）的宇宙显然是一个特例。此时尺度因子按指数而非幂律膨胀；整个度规为
\begin{equation*}
\mathrm{d}s^{2} = -\mathrm{d}t^{2} + e^{Ht}[\mathrm{d}x^{2} + \mathrm{d}y^{2} + \mathrm{d}z^{2}],
\tag{8.89}
\end{equation*}

% ---- 书页 341 / p356 ----
其中 Hubble 参量 $H$ 是常数。当然，在 8.1 节中我们已经描述过一个具有正宇宙学常数的宇宙学时空：de Sitter 空间，其特征是 $\kappa>0$ 且 $a\propto\cosh(t/\alpha)$。那个解与此处 $\kappa=0$、$a\propto\exp(Ht)$ 的解之间是什么关系？它们是同一个时空，只是用不同的坐标来表示。验证这一点的一种办法，是计算 (8.89) 的 Riemann 张量，并检验它具有最大对称时空的特征形式，即式 (8.1)。由于正曲率的最大对称时空在局部是唯一的，度规 (8.6) 与 (8.89) 必定描述同一个流形，或其一部分。实际上，(8.89) 的坐标只覆盖了 de Sitter 空间的一部分；它们在过去是不完备的。在习题中你会被要求证明，这些坐标中的共动测地线会在有限的仿射参量内到达 $t=-\infty$；它们撞上了坐标的边缘。在图 8.1 的共形图中，这些坐标覆盖了正方形右上方的三角形部分。关于 de Sitter 空间和反 de Sitter 空间上各种坐标系的更完备的描述，参见 Hawking and Ellis (1973)。

另一个有趣的特例是 $\rho=0$ 但具有空间曲率的完全空无一物的宇宙。弗里德曼方程变为
\begin{equation*}
H^{2} = -\frac{\kappa}{a^{2}},
\tag{8.90}
\end{equation*}
所以曲率 $\kappa$ 必须为负。把曲率看作一种虚构的能量密度 $\rho_{\mathrm{c}}\propto a^{-2}$，由式 (8.88) 可知，这样的宇宙将线性膨胀，$a\propto t$。这个时空被称为\textbf{Milne 宇宙}（Milne universe）。然而，正如 de Sitter 空间的情形一样，我们知道另一个 $\rho=0$ 的宇宙学时空——在这里就是平直的 Minkowski 空间。再一次，Milne 时空只是 Minkowski 空间在某个不完备坐标系中的一个片。它可以被看作 Minkowski 空间中某个固定点的未来光锥的内部，用负曲率的双曲面来分层。要检验这一点，只需计算 Riemann 张量的所有分量——结果它们全部为零；任何 Riemann 曲率为零的时空在局部都是 Minkowski 空间。

与这些理想化的解相反，一个现实的宇宙学将包含好几种形式的能量-动量。在当前的宇宙中，我们有把握认为辐射密度显著低于物质密度，但真空与物质在动力学上都很重要。因此，方便的做法是用 $\Omega_{\mathrm{M}}$ 和 $\Omega_{\Lambda}$ 来参数化像我们这样的宇宙，曲率则由 $\Omega_{\mathrm{c}}=1-\Omega_{\mathrm{M}}-\Omega_{\Lambda}$ 固定。这类宇宙某些具体例子的膨胀历史示于图 8.3。随着这些宇宙的膨胀，物质、曲率和真空的相对影响会发生变化，因为相应的密度以不同的速率演化：
\begin{equation*}
\Omega_{\Lambda}\propto\Omega_{\mathrm{c}}a^{2}\propto\Omega_{\mathrm{M}}a^{3}.
\tag{8.91}
\end{equation*}
当 $a\to0$ 的过去，曲率和真空将可以忽略，宇宙的行为如同 Einstein--de Sitter。当 $a\to\infty$ 的未来，曲率和物质将可以忽略，宇宙将渐近于 de Sitter 空间；除非尺度因子
% ---- 书页 342 / p357 ----
永远达不到无限大，因为宇宙会在某个有限的时刻开始再度坍缩。

%FIG{8.3}: {不同 $\Omega_{\mathrm{M}}$ 与 $\Omega_{\Lambda}$ 取值下的膨胀历史。自上而下，各条曲线分别描述 $(\Omega_{\mathrm{M}},\Omega_{\Lambda})=(0.3,0.7)$、$(0.3,0.0)$、$(1.0,0.0)$ 和 $(4.0,0.0)$。}

如果真空能为负，就\emph{总会}发生再度坍缩；随着宇宙膨胀，真空能最终将占据主导，而 $\Omega_{\Lambda}<0$ 的效果是导致减速和再度坍缩（正如 $\Omega_{\Lambda}>0$ 的效果是把宇宙推开）。当 $\Omega_{\Lambda}\ge0$ 时，如果 $\Omega_{\mathrm{M}}$ 大到足以在 $\Omega_{\Lambda}$ 来得及接管之前使宇宙的膨胀停下来，再度坍缩也是可能的。这些可能性表现为图 8.4 中 $\Omega_{\mathrm{M}}/\Omega_{\Lambda}$ 参数空间的不同区域。对角线代表 $\Omega_{\mathrm{total}}=1$，意味着 $\kappa=0$。

要确定永恒膨胀与最终再度坍缩之间的分界线，请注意：坍缩要求 Hubble 参量在从正变负的过程中经过零。发生这种转折处的尺度因子 $a_{*}$，可以通过在弗里德曼方程中令 $H=0$ 而求得，
\begin{equation*}
H^{2} = 0 = \frac{8\pi G}{3}\left(\rho_{\mathrm{M}0}a_{*}^{-3} + \rho_{\Lambda0} + \rho_{\mathrm{c}0}a_{*}^{-2}\right).
\tag{8.92}
\end{equation*}
把它除以 $H_0^{2}$，利用 $\Omega_{\mathrm{c}0}=1-\Omega_{\mathrm{M}0}-\Omega_{\Lambda0}$，再稍加整理，便得到
\begin{equation*}
\Omega_{\Lambda0}a_{*}^{3} + (1-\Omega_{\mathrm{M}0}-\Omega_{\Lambda})a_{*} + \Omega_{\mathrm{M}0} = 0.
\tag{8.93}
\end{equation*}
（译注：按式 (8.92) 的推导，第二项应为 $\Omega_{\Lambda0}$；原书排印漏写下标 0，原文如此。）

这是关于 $a_{*}$（转折点处的尺度因子）的三次方程。当然，我们对 $a_{*}$ 本身并不很在乎；我们在乎的是，给定 $\Omega_{\mathrm{M}0}$，使 (8.93) 存在实解的 $\Omega_{\Lambda0}$ 值。求解这个三次方程并做一些数学运算，我们发现，宇宙将永远膨胀所对应的 $\Omega_{\Lambda0}$ 值由下式给出

% ---- 书页 343 / p358 ----
%FIG{8.4}: {由物质和真空能主导的宇宙的性质，表示为密度参数 $\Omega_{\mathrm{M}}$ 与 $\Omega_{\Lambda}$ 的函数。左上角的圆形区域大致代表实验数据（截至 2003 年）所偏好的数值。}

\begin{equation*}
\Omega_{\Lambda0} \ge
\left\{
\begin{array}{ll}
0 & 0\le\Omega_{\mathrm{M}0}\le1 \\
4\Omega_{\mathrm{M}0}\cos^{3}\left[\dfrac{1}{3}\cos^{-1}\left(\dfrac{1-\Omega_{\mathrm{M}0}}{\Omega_{\mathrm{M}0}}\right) + \dfrac{4\pi}{3}\right] & \Omega_{\mathrm{M}0}>1.
\end{array}
\right.
\tag{8.94}
\end{equation*}
注意，当 $\Omega_{\Lambda0}=0$ 时，开放和平直宇宙（$\Omega_0=\Omega_{\mathrm{M}0}\le1$）将永远膨胀，而闭合宇宙（$\Omega_0=\Omega_{\mathrm{M}0}>1$）将再度坍缩。对宇宙学常数由来已久的轻视，造成了一种民间信念，认为这是一种必然的对应关系；然而，一旦承认真空能的可能性，空间几何与最终命运之间的任何组合就都是可能的了。

在图 8.4 的左上角，我们标出了目前受偏好的宇宙学参数取值：$\Omega_{\mathrm{M}0}\sim0.3$、$\Omega_{\Lambda0}\sim0.7$，我们将在 8.7 节中讨论。这已深入永恒膨胀的区域；如果真空能保持真正的恒定（它未必如此），我们的宇宙就注定要永远膨胀下去。

在本节的结尾，我们指出为弗里德曼方程寻找静态解的困难。要成为静态，我们必须不仅有 $\dot{a}=0$，还要 $\ddot{a}=0$。由式 (8.68)，这只有在压强为
\begin{equation*}
p = -\frac{1}{3}\rho,
\tag{8.95}
\end{equation*}
时才可能；而由式 (8.67)，还必须存在非零的空间曲率
\begin{equation*}
\frac{\kappa}{a^{2}} = \frac{8\pi G}{3}\rho.
\tag{8.96}
\end{equation*}

% ---- 书页 344 / p359 ----
由于能量密度和压强必须异号，如果我们只诉诸物质或辐射，这些条件就无法满足。当 Einstein 最初在广义相对论中寻找宇宙学解时，天文学家尚未发现宇宙正在膨胀，因此静态解的缺失曾被认为是个难题。这为 Einstein 引入宇宙学常数提供了动机；静态条件可以由物质和真空能的组合来满足，其中
\begin{equation*}
\rho_{\Lambda} = \frac{1}{2}\rho_{\mathrm{M}},
\tag{8.97}
\end{equation*}
再加上适当的正空间曲率。这些参数描述的是\textbf{Einstein 静态宇宙}（Einstein static universe）。如今我们知道宇宙在膨胀，因此这个解没有多少经验上的意义；但它对理论家却极为有用，为共形图的构造提供了基础。

\section{红移与距离}

显然，我们希望通过观测测定若干个量，以判断 FRW 模型中的哪一个对应我们的宇宙。显然我们想测定 $H_0$，因为它与宇宙的年龄有关；我们还想知道 $\Omega$，它通过式 (8.76) 决定 $\kappa$。为了理解这些量究竟可能如何被测量，让我们考虑 FRW 宇宙中的测地线运动。FRW 宇宙中存在若干类空 Killing 矢量，却没有能给我们守恒能量概念的类时 Killing 矢量。不过，确实存在一个 Killing 张量。如果 $U^{\mu}=(1,0,0,0)$ 是共动观者的四速度，那么张量
\begin{equation*}
K_{\mu\nu} = a^{2}(g_{\mu\nu} + U_{\mu}U_{\nu})
\tag{8.98}
\end{equation*}
满足 $\nabla_{(\sigma}K_{\mu\nu)}=0$（你可以自行验证），因而是一个 Killing 张量。这意味着，如果粒子具有四速度 $V^{\mu}=\mathrm{d}x^{\mu}/\mathrm{d}\lambda$，那么
\begin{equation*}
K^{2} = K_{\mu\nu}V^{\mu}V^{\nu} = a^{2}[V_{\mu}V^{\mu} + (U_{\mu}V^{\mu})^{2}]
\tag{8.99}
\end{equation*}
沿测地线将是常数。让我们先就有质量粒子的情形来思考这个问题。此时我们将有 $V_{\mu}V^{\mu}=-1$，于是
\begin{equation*}
(V^{0})^{2} = 1 + |\vec{V}|^{2},
\tag{8.100}
\end{equation*}
其中 $|\vec{V}|^{2}=g_{ij}V^{i}V^{j}$。我们还有 $U_{\mu}V^{\mu}=-V^{0}$，所以式 (8.99) 给出
\begin{equation*}
|\vec{V}| = \frac{K}{a}.
\tag{8.101}
\end{equation*}
因此，随着宇宙膨胀，粒子相对于共动坐标会“慢下来”。事实上这是一种真实的变慢，其意义是：初始相对速度很高的一团粒子气体将随宇宙膨胀而冷却下来。

% ---- 书页 345 / p360 ----
类似的事情也发生在类光测地线上。此时 $V_{\mu}V^{\mu}=0$，式 (8.99) 给出
\begin{equation*}
U_{\mu}V^{\mu} = \frac{K}{a}.
\tag{8.102}
\end{equation*}
但共动观者测得的光子频率是 $\omega=-U_{\mu}V^{\mu}$。因此，以频率 $\omega_{\mathrm{em}}$ 发出的光子，随着宇宙的膨胀将被观测到较低的频率 $\omega_{\mathrm{obs}}$：
\begin{equation*}
\frac{\omega_{\mathrm{obs}}}{\omega_{\mathrm{em}}} = \frac{a_{\mathrm{em}}}{a_{\mathrm{obs}}}.
\tag{8.103}
\end{equation*}
宇宙学家喜欢用这两个事件之间的\textbf{红移}（redshift）$z$ 来谈论它，红移由波长的分数变化定义：
\begin{equation*}
z_{\mathrm{em}} = \frac{\lambda_{\mathrm{obs}} - \lambda_{\mathrm{em}}}{\lambda_{\mathrm{em}}}.
\tag{8.104}
\end{equation*}
如果观测发生在今天（$a_{\mathrm{obs}}=a_0=1$），这意味着
\begin{equation*}
\boxed{\,a_{\mathrm{em}} = \frac{1}{1+z_{\mathrm{em}}}.\,}
\tag{8.105}
\end{equation*}
所以，一个天体的红移告诉我们光子发射时的尺度因子。

注意，这种红移与传统的多普勒效应并不是一回事；导致红移的是空间的膨胀，而不是观者与发射者的相对速度。不过，如果我们在与 Hubble 半径 $H_0^{-1}$ 以及空间曲率半径 $\kappa^{-1/2}$ 相比很小的距离上观测星系，宇宙的膨胀看上去就很像一群星系在彼此远离，红移看上去也很像多普勒效应。因此，天文学家常常用“速度”$v=cz$（$c$ 是光速）来考虑红移。尽管我们知道，对弯曲时空中不同点上的两个物体并不能真正谈论相对速度，但这个虚构的说法在足够短的距离上工作得很好。在这个近似之内，从我们到某个星系的“距离”$d$ 可以取为\textbf{瞬时物理距离}（instantaneous physical distance）$d_P$（即在诸如厘米这样的物理单位下，沿我们当前的空间超曲面量出的、我们与该星系位置之间的距离）。让我们把 RW 度规写成
\begin{equation*}
\mathrm{d}s^{2} = -\mathrm{d}t^{2} + a^{2}(t)R_0^{2}\left[\mathrm{d}\chi^{2} + S_k^{2}(\chi)\,\mathrm{d}\Omega^{2}\right],
\tag{8.106}
\end{equation*}
其中 $S_k(\chi)$ 由式 (8.33) 定义，而 $k\in\{+1,0,-1\}$。在这种形式下，在时刻 $t$ 测得的、我们（$\chi=0$）与共动径向坐标为 $\chi$ 的星系之间的瞬时物理距离就是
\begin{equation*}
d_P(t) = a(t)R_0\chi,
\tag{8.107}
\end{equation*}

% ---- 书页 346 / p361 ----
其中 $\chi$ 保持不变，因为我们假定我们和被观测的星系都完美地共动。（它们可能并不如此；在那种情形下，把所谓“本动速度”（peculiar velocities）引起的修正包括进来是轻而易举的。）当然，“距离”之所以加引号，是因为一旦离开这个近似，就有好几种彼此不等价而有用的距离概念；但当 $d_P$ 很小时它们全都一致。这时，观测到的速度（由红移推断）就简单地是
\begin{equation*}
v = \dot{d}_P = \dot{a}R_0\chi = \frac{\dot{a}}{a}d_P.
\tag{8.108}
\end{equation*}
在今天取值，它就成为
\begin{equation*}
\boxed{\,v = H_0 d_P,\,}
\tag{8.109}
\end{equation*}
这就是著名的\textbf{Hubble 定律}（Hubble law）：对于不太遥远的星系，观测到的退行速度与距离成正比。

如果红移不是非常小，我们就得更仔细地思考宇宙学中“距离”指的是什么。瞬时物理距离是一个方便的构造，但它本身不可观测，因为观测涉及的总是我们过去光锥上的事件，而不是我们当前的空间超曲面。在欧几里得空间中，推断一个物体的距离有许多不同的方法；例如，我们可以把它的视亮度与它的内禀光度相比较，或者把它的视角速度与它的内禀横向速度相比较，或者把它的视角大小与它的物理尺度相比较。对每一种情形，我们都可以定义一种距离，即假如空间是欧几里得的、宇宙不在膨胀时我们\emph{将会}推断出的距离。

让我们从\textbf{光度距离}（luminosity distance）$d_L$ 开始，它的定义是满足
\begin{equation*}
d_L^{2} = \frac{L}{4\pi F},
\tag{8.110}
\end{equation*}
其中 $L$ 是源的绝对光度，$F$ 是观者测得的流量（flux），即单位时间落到某个探测器单位面积上的能量。这个定义来自如下事实：在平直空间中，对位于距离 $d$ 处的源，流量与光度之比正好是以源为中心的球面面积的倒数，$F/L = 1/A(d) = 1/4\pi d^{2}$。然而在 FRW 宇宙中，流量将被稀释。光子数守恒告诉我们，源发出的所有光子最终都将穿过与发射者相距共动距离 $\chi$ 的球面。但流量还会被两个附加效应稀释：单个光子红移一个因子 $(1+z)$；而且光子撞击球面的频率变低，因为相隔 $\delta t$ 发出的两个光子，被观测到的相隔时间是 $(1+z)\delta t$。于是我们将有
\begin{equation*}
\frac{F}{L} = \frac{1}{(1+z)^{2}A}.
\tag{8.111}
\end{equation*}

% ---- 书页 347 / p362 ----
以共动距离 $\chi$ 为中心的球面的面积 $A$，可以从式 (8.106) 中 $\mathrm{d}\Omega^{2}$ 的系数导出，得到
\begin{equation*}
A = 4\pi R_0^{2}S_k^{2}(\chi),
\tag{8.112}
\end{equation*}
其中我们已取 $a(t)=1$，因为我们是在今天观测这些光子。把所有这些合在一起，就得到
\begin{equation*}
d_L = (1+z)R_0S_k(\chi).
\tag{8.113}
\end{equation*}
光度距离 $d_L$ 是我们有希望测量的量，因为有一些天体物理源的绝对光度是已知的。但 $\chi$ 是不可观测的，所以我们必须把它从方程中消去。在类光测地线（为方便起见取为径向）上，我们有
\begin{equation*}
0 = \mathrm{d}s^{2} = -\mathrm{d}t^{2} + a^{2}R_0^{2}\mathrm{d}\chi^{2},
\tag{8.114}
\end{equation*}
即
\begin{equation*}
\chi = R_0^{-1}\int\frac{\mathrm{d}t}{a} = R_0^{-1}\int\frac{\mathrm{d}a}{a^{2}H(a)},
\tag{8.115}
\end{equation*}
其中我们用了 $H=\dot{a}/a$。惯常的做法是用 $a=1/(1+z)$ 把尺度因子换成红移，于是我们有
\begin{equation*}
\chi(z) = R_0^{-1}\int_{0}^{z}\frac{\mathrm{d}z'}{H(z')}.
\tag{8.116}
\end{equation*}

为了计算这个积分中的 Hubble 参量，我们使用弗里德曼方程 (8.67)，像上一节中那样把它写成
\begin{equation*}
H^{2} = \frac{8\pi G}{3}\sum_{i(\mathrm{c})}\rho_i.
\tag{8.117}
\end{equation*}
为了简化问题，我们可以再次假设每个密度组分都按幂律演化，
\begin{equation*}
\rho_i(z) = \rho_{i0}a^{-n_i} = \rho_{i0}(1+z)^{n_i},
\tag{8.118}
\end{equation*}
于是我们可以写出
\begin{equation*}
H(z) = H_0E(z),
\tag{8.119}
\end{equation*}
其中
\begin{equation*}
E(z) = \left[\sum_{i(\mathrm{c})}\Omega_{i0}(1+z)^{n_i}\right]^{1/2},
\tag{8.120}
\end{equation*}

% ---- 书页 348 / p363 ----
其中密度参数 $\Omega_i$ 由式 (8.74) 定义。下面涉及 $E(z)$ 的那些方程，无论能量源是否按幂律演化都成立；若不按幂律演化，只需采用 $E(z)=H(z)/H_0$［其中 $H(z)$ 由弗里德曼方程确定］而不必用式 (8.120)。

于是，光度距离为
\begin{equation*}
d_L(z) = (1+z)R_0S_k\left[R_0^{-1}H_0^{-1}\int\frac{\mathrm{d}z'}{E(z')}\right].
\tag{8.121}
\end{equation*}
注意，当 $k=0$ 时 $R_0$ 自动消去，这很好，因为在那种情形下它是一个完全任意的参数。即使它并非任意，人们通常也更愿意用 $\Omega_{\mathrm{c}0}=-k/R_0^{2}H_0^{2}$ 来讨论；这个量既可以由对空间曲率的测定直接测量，也可以通过测量密度参数并利用 $\Omega_{\mathrm{c}0}=1-\Omega_0$ 来得到。用这个参数，我们有
\begin{equation*}
R_0 = H_0^{-1}\sqrt{-k\Omega_{\mathrm{c}0}} = \frac{H_0^{-1}}{\sqrt{|\Omega_{\mathrm{c}0}|}}.
\tag{8.122}
\end{equation*}
于是，我们把光度距离用可测量的宇宙学参数写成
\begin{equation*}
\boxed{\,d_L(z) = (1+z)\frac{H_0^{-1}}{\sqrt{|\Omega_{\mathrm{c}0}|}}S_k\left[\sqrt{|\Omega_{\mathrm{c}0}|}\int\frac{\mathrm{d}z'}{E(z')}\right].\,}
\tag{8.123}
\end{equation*}
这个方程虽然显得笨拙，却在宇宙学中处于核心地位。给定可观测量 $H_0$ 和 $\Omega_{i0}$，我们可以直接算出到位于任一红移 $z$ 处天体的光度距离；同样地，我们也可以对一系列红移范围内的天体测量 $d_L(z)$，并从这些信息中提取 $H_0$ 和/或各个 $\Omega_{i0}$。

除光度距离之外，还有两个相关的距离量度。正如光度距离是假如我们身处平直空间时，从源的内禀光度和观测光度推断出的距离，\textbf{固有运动距离}（proper motion distance）$d_M$ 是从源的内禀运动和观测运动推断出的距离。它的定义为
\begin{equation*}
d_M = \frac{u}{\dot{\theta}},
\tag{8.124}
\end{equation*}
其中 $u$ 是固有横向速度（比如你会以 k/s 为单位来测量的量），而 $\dot{\theta}$ 是观测到的角速度。另一方面，\textbf{角直径距离}（angular diameter distance）是从源的内禀大小和观测大小推断出的距离；它的定义为
\begin{equation*}
d_A = \frac{R}{\theta},
\tag{8.125}
\end{equation*}
